\chapter{Introduction}\label{ch:introduction}

\section{What WOP does}
WOP (Workshop Optimised Programming) makes G-code for a CNC lathe. You do not write G-code. You describe the job: the raw part, the shape you want, the tools and what to do in which order. WOP works out the tool paths, shows them in a picture, checks them, and writes the G-code for your controller.

WOP is made for lathes with two axes (X and Z) that run on LinuxCNC. Other controllers are possible through the postprocessors (Chapter~\ref{ch:posts}), but only LinuxCNC is the target.

\section{How a program is made}
A program is a list of \emph{blocks}. Each block is one job on the part: facing, turning, drilling, threading, grooving and so on. The blocks are done from top to bottom. A block works on the material that the blocks above it left over.

You work in this order:
\begin{enumerate}
  \item Set the machine, the raw part and the workholding in the \emph{Setup} (Chapter~\ref{ch:project}).
  \item Add tools (Chapter~\ref{ch:tools}).
  \item Add blocks and fill in their forms (Chapter~\ref{ch:blocks}). A block that follows a shape needs a \emph{contour}, which you draw in the sketch editor (Chapter~\ref{ch:contours}).
  \item Look at the picture, step through the moves and read the messages (Chapter~\ref{ch:checking}).
  \item Export the G-code (Chapter~\ref{ch:file}, Section~\ref{sec:file:export}).
\end{enumerate}
You can go back at any step. WOP plans the whole program again after every change.

\section{The words used in this manual}
\subsection{Project}
A project is one job in one file (\file{.wop.json}). It holds the Setup, the blocks, the contours, and the tools and workholding of the job.

\subsection{Block}
A block is one step of the program. The blocks are listed in the flow list on the left of the window. A block has a form on the right with its values.

\subsection{Contour}
A contour is a shape drawn in the sketch editor, a line of straight lines and arcs. A block can follow a contour. A contour is listed under the heading \emph{Contours} in the flow list. It is not a step of the program, and it can be copied.

\subsection{Tool and workholding}
A \emph{tool} is a cutting tool with its real shape. \emph{Workholding} is the chuck body and the jaws that hold the part. WOP uses their shapes to find collisions.

\subsection{Machine and postprocessor}
A \emph{machine} holds what belongs to your lathe: the travel of the axes, the spindle speed and power, the tool change position and which postprocessor to use. A \emph{postprocessor} (post) is a text file that says how the G-code is written for a controller.

\section{Coordinates and units}
WOP uses millimetres.

\subsection{X as radius or diameter}
X is the distance from the centre line of the part. In the project you choose whether sizes in X are typed and shown as a \emph{radius} or as a \emph{diameter} (Section~\ref{sec:project:xmode}). The labels of the fields say which. The G-code has its own setting in the postprocessor.

\subsection{Z and the picture}
Z runs along the part. The picture shows Z to the right and X up; the chuck is on the left. The face of the part is at $Z=0$ and the part lies towards minus Z.

\section{Where WOP is limited}
\begin{itemize}
  \item WOP has two axes. There is no C axis, no milling and no live tools.
  \item The G-code is written as drawn. WOP does not shorten a path because a tool might not reach; the check reports it and you decide (Chapter~\ref{ch:checking}).
  \item Nothing is tried on every machine. Parts that are not confirmed on a real machine are marked in this manual. Check what a postprocessor writes before you run it; any G-code can crash a machine.
  \item The checks are estimates. The power check is an approximation, and the check of the feed speed is not a full simulation of the controller.
\end{itemize}
